% =========================================================
% PENJELASAN DETAIL: Unit Commitment dengan Virtual Inertia
% Author  : Nathaniel Benelisha (UGM)
% Version : v1.0 - September 2026
% Berdasarkan: Coba_16_20260612_R00_vinertia_cplexdirect.ipynb
%              main.tex, draft_paper.md, dan Referensi/
% =========================================================

\documentclass[journal]{IEEEtran}
\usepackage{amsmath,amssymb}
\usepackage{booktabs}
\usepackage{cite}
\usepackage{xcolor}
\usepackage{listings}
\usepackage{array}

\lstdefinestyle{pycode}{
  language=Python,
  basicstyle=\ttfamily\footnotesize,
  keywordstyle=\color{blue},
  commentstyle=\color{gray},
  numbers=left,
  numbersep=5pt,
  frame=single,
  breaklines=true,
  captionpos=b,
  showstringspaces=false,
  tabsize=4
}

\newcommand{\Hsys}{H_{\mathrm{sys}}}
\newcommand{\Ploss}{\Delta P_{\mathrm{loss}}}
\newcommand{\RoCoF}{\mathrm{RoCoF}}
\newcommand{\HVIBESS}{H_{\mathrm{VI,BESS}}}
\newcommand{\HVIWT}{H_{\mathrm{VI,WT}}}
\newcommand{\HVIPV}{H_{\mathrm{VI,PV}}}

\begin{document}

\title{Penjelasan Detail: Unit Commitment dengan Virtual Inertia\\
dari Sumber Daya Berbasis Inverter (IBR)\\
--- Implementasi dan Teori ---}

\author{Nathaniel~Benelisha,~Universitas~Gadjah~Mada\\
\texttt{nathanielbenelisha@mail.ugm.ac.id}}

\maketitle

\begin{abstract}
Dokumen ini menyajikan penjelasan menyeluruh mengenai formulasi
\textit{Frequency-Constrained Unit Commitment} (FCUC) yang
mempertimbangkan kontribusi \textit{virtual inertia} (VI)
dari berbagai sumber daya berbasis inverter (IBR), meliputi
BESS (\textit{Battery Energy Storage System}), turbin angin (WT),
dan solar PV. Penjelasan mencakup: (1) landasan fisik inersia sistem
dan respons frekuensi pasca-gangguan; (2) model matematis setiap
komponen VI; (3) formulasi lengkap MILP FCUC; dan (4) korespondensi
formulasi matematis dengan implementasi Python/Pyomo pada
\texttt{Coba\_16\_20260612\_R00\_vinertia\_cplexdirect.ipynb}.
\end{abstract}

\begin{IEEEkeywords}
Virtual Inertia, FCUC, IBR, BESS, Wind Turbine, Solar PV,
RoCoF, Frequency Nadir, MILP, Pyomo, CPLEX.
\end{IEEEkeywords}

%-----------------------------------------------------------
\section{Pendahuluan dan Motivasi}
\label{sec:intro}
%-----------------------------------------------------------

\subsection{Masalah Inersia dalam Sistem Tenaga Modern}

Sistem tenaga listrik konvensional mengandalkan inersia
rotasional dari generator sinkron (SG) sebagai mekanisme
alami untuk menahan laju perubahan frekuensi
(\textit{Rate of Change of Frequency}, RoCoF) setelah
terjadi gangguan kehilangan unit pembangkit (kontigensi N-1).

Secara fisik, inersia rotasional dimodelkan melalui
\textit{swing equation}:
\begin{equation}
  2H_i \frac{d\Delta f}{dt} = \frac{P_m - P_e}{S_i}
  \label{eq:swing_basic}
\end{equation}
di mana $H_i$ [MWs/MVA] adalah konstanta inersia unit $i$,
$P_m$ daya mekanik, $P_e$ daya listrik, dan $S_i$ kapasitas
terpasang. Untuk sistem agregat dengan $N_g$ generator:
\begin{equation}
  2\Hsys \frac{d\Delta f}{dt} =
    -\Ploss + \Delta P_{\mathrm{gov}}(t)
    + \Delta P_{\mathrm{VI}}(t) - D\,\Delta f(t)
  \label{eq:swing_sys}
\end{equation}
di mana $\Hsys$ adalah inersia sistem total [MWs/MVA],
$\Ploss$ adalah daya yang hilang, $\Delta P_{\mathrm{gov}}$
adalah respons governor, $\Delta P_{\mathrm{VI}}$ adalah
injeksi daya virtual inertia, dan $D$ adalah koefisien
redaman beban.

\subsection{Dampak Tingginya Penetrasi IBR}

IBR --- termasuk PV, turbin angin, dan BESS --- tidak
memiliki inersia rotasional secara alami. Akibatnya:
\begin{itemize}
  \item \textbf{RoCoF meningkat}: Frekuensi jatuh lebih
    cepat, berisiko memicu relay proteksi~\cite{Tielens2016}.
  \item \textbf{Nadir frekuensi lebih dalam}: Titik terendah
    frekuensi dapat turun di bawah ambang UFLS
    (49{,}0~Hz untuk sistem Indonesia)~\cite{ESDM2007}.
  \item \textbf{QSS memburuk}: Frekuensi keadaan tunak
    setelah respons primer dapat tidak memenuhi
    persyaratan minimum.
\end{itemize}

RUPTL PLN 2025--2034 menargetkan 61\% kapasitas dari
energi terbarukan pada 2034~\cite{PLN2025}, sehingga
pemodelan inersia virtual menjadi kritis untuk perencanaan
operasi sistem.

\subsection{Konsep Virtual Inertia}

\textit{Virtual inertia} (VI) adalah kemampuan IBR
\textbf{meniru} perilaku inersia rotasional melalui
algoritma kontrol konverter daya:
\begin{enumerate}
  \item \textbf{Grid-forming converter (GFM)}: Konverter
    beroperasi dengan referensi tegangan, merespons
    $d\Delta f/dt$ dengan memodulasi daya~\cite{Asiedu2025}.
  \item \textbf{Synthetic inertia loop} pada turbin angin:
    Energi kinetik rotor diekstraksi sementara menggunakan
    algoritma dua tahap~\cite{Chu2020}.
  \item \textbf{Fast power injection} dari BESS: BESS
    menyuntikkan daya dalam milidetik melalui loop
    kontrol VI~\cite{Wen2016}.
\end{enumerate}

%-----------------------------------------------------------
\section{Pemodelan Komponen Sistem}
\label{sec:model}
%-----------------------------------------------------------

\subsection{Generator Sinkron Konvensional}

Biaya bahan bakar unit $i$ pada periode $t$:
\begin{equation}
  C_i^{\mathrm{gen}}(P_{i,t}) = a_i P_{i,t}^2 + b_i P_{i,t} + c_i
  \label{eq:fuel_cost}
\end{equation}
Untuk menjaga linieritas MILP, digunakan aproksimasi
\textit{piecewise linear} (PWL):
\begin{equation}
  C_i^{\mathrm{gen}} \approx
    c_i\,u_{i,t} + \sum_{k=1}^{K} \lambda_{i,k}\,\Delta P_{i,k,t}
  \label{eq:pwl}
\end{equation}
Kontribusi inersia setiap unit thermal yang beroperasi:
$H_{\mathrm{SG},i} \cdot S_i \cdot u_{i,t}$ [MWs],
dengan $H_{\mathrm{SG},i}$ tipikal 2--9 MWs/MVA~\cite{Kundur1994}.

\textit{Implementasi (Cell 12):} Parameter $H_{\mathrm{SG},i}$
disimpan di kolom \texttt{``H''} pada file Excel, dimuat ke
parameter Pyomo \texttt{model.H[g]}.

\subsection{Model BESS}

Dinamika \textit{State of Charge} (SoC):
\begin{equation}
  E_{b,t} = E_{b,t-1}
    + \eta_c^b \, P_{c,b,t}\,\Delta t
    - \frac{1}{\eta_d^b} P_{d,b,t}\,\Delta t
  \label{eq:soc}
\end{equation}
di mana $E_{b,t}$ [MWh] adalah energi tersimpan, $\eta_c^b$
dan $\eta_d^b$ adalah efisiensi, $P_{c,b,t}$ dan $P_{d,b,t}$
[MW] adalah daya pengisian dan pengosongan.

\begin{lstlisting}[style=pycode,caption={Dinamika SoC BESS (Cell 52)}]
def battery_soc_rule(m, b, t):
    pos = T_list.index(t)
    if pos == 0:
        return (m.SOC[b, t] == m.SOC_init[b]
                + m.eta_ch[b]*m.P_charge[b, t]
                - (1/m.eta_dis[b])*m.P_discharge[b, t])
    t_prev = T_list[pos - 1]
    return (m.SOC[b, t] == m.SOC[b, t_prev]
            + m.eta_ch[b]*m.P_charge[b, t]
            - (1/m.eta_dis[b])*m.P_discharge[b, t])
\end{lstlisting}

\subsubsection{Virtual Inertia dari BESS}

BESS menyediakan VI melalui kontrol \textit{grid-forming}:
\begin{equation}
  \HVIBESS(t) = K_{\mathrm{VI}}^{\mathrm{BESS}} \cdot P_{\mathrm{VI},b,t}
  \label{eq:h_bess_vi}
\end{equation}
di mana $K_{\mathrm{VI}}^{\mathrm{BESS}}$ [MWs/MW] adalah
gain inersia virtual. Constraint discharge dengan headroom VI:
\begin{equation}
  P_{d,b,t} + \frac{P_{\mathrm{VI},b,t}}{\eta_{\mathrm{VI}}^b}
    \leq P_{d,b}^{\max} \cdot \delta_{d,b,t}
  \label{eq:bess_dis_vi}
\end{equation}

\begin{lstlisting}[style=pycode,caption={Constraint discharge BESS dengan VI headroom (Cell 52)}]
def battery_discharge_limit_rule(m, b, t):
    if CONFIG["con_battery_vi"]:
        return (
            m.P_discharge[b, t]
            + m.P_VI_batt[b, t] / m.Eff_VI[b]
            <= m.DR_max[b] * m.DischargeStatus[b, t]
        )
    return (m.P_discharge[b, t]
            <= m.DR_max[b] * m.DischargeStatus[b, t])
\end{lstlisting}

\textbf{Catatan:} $P_{\mathrm{VI},b,t}$ adalah cadangan
yang dijadwalkan (\textit{ex-ante}), bukan injeksi real-time.
Model memastikan \textit{headroom} tersedia; konverter
menggunakannya otomatis saat gangguan~\cite{Wen2016}.

\subsection{Model Turbin Angin (WT)}

Inersia sintetis dari turbin angin melalui ekstraksi energi
kinetik rotor~\cite{Chu2020}:
\begin{equation}
  \HVIWT(t) = K_{\mathrm{VI}}^{\mathrm{WT}} \cdot P_{\mathrm{VI,WT},w,t}
  \label{eq:h_wt_vi}
\end{equation}

\textbf{Mekanisme fisik:} Rotor turbin memiliki energi
kinetik $E_k = \frac{1}{2}J\omega^2$. Saat frekuensi turun,
konverter meningkatkan ekstraksi daya sementara (melebihi
MPPT), rotor melambat melepas $E_k$ ke grid. Kontrol
dua-tahap dari~\cite{Chu2020} menangani pemulihan rotor
untuk menghindari \textit{secondary frequency dip}.

VI dari WT hanya diambil dari porsi yang dikurtail:
\begin{equation}
  P_{\mathrm{VI,WT},w,t} \leq P_{\mathrm{Curt,WT},w,t}
  \label{eq:wt_vi_curtailment}
\end{equation}

\begin{lstlisting}[style=pycode,caption={Constraint VI WT dari curtailment (Cell 55)}]
if CONFIG["con_wt"] and CONFIG["con_wt_vi"]:
    def wt_vi_curtailment_rule(m, w, t):
        return m.P_VI_wt[w, t] <= m.WT_Curt[w, t]
    model.WTVICurtailmentLimit = Constraint(
        model.WT, model.T, rule=wt_vi_curtailment_rule)
\end{lstlisting}

\subsection{Model Solar PV}

VI dari PV melalui \textit{grid-forming inverter} dan
\textit{deloading}~\cite{Asiedu2025}:
\begin{equation}
  \HVIPV(t) = K_{\mathrm{VI}}^{\mathrm{PV}} \cdot P_{\mathrm{VI,PV},t}
  \label{eq:h_pv_vi}
\end{equation}
Dalam kode \texttt{Coba\_16}, VI dari PV dimodelkan
serupa dengan VI dari WT.

\subsection{Inersia Sistem Total}

\begin{equation}
  \Hsys(t) =
    \underbrace{\sum_{i \in \mathcal{G}_{\mathrm{th}}}
      H_{\mathrm{SG},i}\,S_i\,u_{i,t}}_{\text{Inersia rotasional (SG)}}
    +
    \underbrace{K_{\mathrm{VI}}^{\mathrm{BESS}} P_{\mathrm{VI,BESS},t}
    + K_{\mathrm{VI}}^{\mathrm{WT}} P_{\mathrm{VI,WT},t}}_{\text{Virtual inertia (IBR)}}
  \label{eq:hsys}
\end{equation}

\begin{lstlisting}[style=pycode,caption={Ekspresi Hsys di Pyomo (Cell 22)}]
def hsys_rule(m, t):
    h_thermal = sum(
        m.H[g] * m.Pmax[g] * m.u[g, t]
        for g in m.G if Type[g] not in ["PV", "WT"]
    )
    h_batt_vi = sum(
        m.Kb_VI[b] * m.P_VI_batt[b, t] for b in m.BATT
    ) if CONFIG["con_battery"] and CONFIG["con_battery_vi"] else 0
    h_wt_vi = sum(
        m.Kw_VI[w] * m.P_VI_wt[w, t] for w in m.WT
    ) if CONFIG["con_wt"] and CONFIG["con_wt_vi"] else 0
    return h_thermal + h_batt_vi + h_wt_vi
model.Hsys = Expression(model.T, rule=hsys_rule)
\end{lstlisting}

\textit{Catatan implementasi:} Pada kode, inersia SG
dihitung sebagai $H[g] \times P_{\max}[g] \times u[g,t]$
karena data $H[g]$ dinormalisasi terhadap $P_{\max}$
sebagai \textit{proxy} MVA.

%-----------------------------------------------------------
\section{Model Respons Frekuensi Sistem}
\label{sec:sfr}
%-----------------------------------------------------------

Model SFR mengikuti Anderson dan Mirheydar~\cite{Anderson1990}.

\subsection{Rate of Change of Frequency (RoCoF)}

\begin{equation}
  \left.\frac{d(\Delta f)}{dt}\right|_{t=0^+}
    = -\frac{f_0 \cdot \Ploss}{2\,\Hsys(t)\,S_{\mathrm{sys}}}
  \label{eq:rocof}
\end{equation}
Batas RoCoF: $|\RoCoF| \leq \RoCoF^{\max}$ (0{,}5 Hz/s
tipikal; 0{,}55 Hz/s dalam kode).

\textbf{Justifikasi VI:} Dengan VI aktif sejak $t=0^+$,
VI meningkatkan $\Hsys$ efektif dan mengurangi RoCoF.
Inilah mengapa inersia virtual masuk langsung ke $\Hsys$.

\begin{lstlisting}[style=pycode,caption={Estimasi RoCoF (Cell 25)}]
def rocof_est_rule(m, t):
    return m.LargestLoss[t] * f0 / (2 * m.Hsys[t])
model.RoCoF_est = Expression(model.T, rule=rocof_est_rule)
\end{lstlisting}

\subsection{Nadir Frekuensi (Linearisasi Dua-Tahap)}

Nadir dilinearisasi menggunakan pendekatan dari~\cite{Shi2025}:
\begin{equation}
  \Delta f_{\mathrm{nadir},t} \approx
    \alpha_0 + \alpha_1 \Hsys(t) + \alpha_2 R_{\mathrm{sys},t}
    + \alpha_3 \Ploss(t)
  \label{eq:nadir_linear}
\end{equation}
Koefisien $\alpha_0,\ldots,\alpha_3$ diperoleh dari regresi
\textit{least-squares} atas 5000 sampel simulasi SFR.

\subsection{Frekuensi Quasi-Steady-State (QSS)}

\begin{equation}
  \Delta f_{\mathrm{QSS}} =
    \frac{f_0 \cdot \Ploss}{D\,S_{\mathrm{sys}} + R_{\mathrm{sys},t}}
  \label{eq:qss}
\end{equation}
di mana $R_{\mathrm{sys},t} = \sum_{i} R_i\,u_{i,t}$
adalah droop governor agregat [MW/Hz].

\begin{lstlisting}[style=pycode,caption={Estimasi frekuensi QSS (Cell 26)}]
def qss_freq_rule(m, t):
    return f0 - (
        m.LargestLoss[t]
        / (m.TotalPFR[t] + D_frac * model.Demand[t])
    )
model.f_qss = Expression(model.T, rule=qss_freq_rule)
\end{lstlisting}

%-----------------------------------------------------------
\section{Formulasi FCUC dengan Virtual Inertia}
\label{sec:formulation}
%-----------------------------------------------------------

\subsection{Fungsi Objektif}

\begin{equation}
  \min \sum_{t=1}^{T} \Bigg[
    \sum_{i \in \mathcal{G}} \Big(
      C_i^{\mathrm{gen}}(P_{i,t})
      + C_i^{\mathrm{SU}} v_{i,t}
      + C_i^{\mathrm{SD}} w_{i,t}
      + c_i^{\mathrm{fixed}} u_{i,t}
    \Big)
    + c_{\mathrm{deg}}\,(P_{c,t}+P_{d,t})\,\Delta t
  \Bigg]
  \label{eq:obj}
\end{equation}

\begin{lstlisting}[style=pycode,caption={Fungsi objektif modular (Cell 28)}]
def obj_rule(m):
    total_cost = 0
    if CONFIG["obj_generation_cost"]:
        total_cost += sum(m.b[g]*m.P[g,t]
                         for g in m.G for t in m.T)
    if CONFIG["obj_fixed_cost"]:
        total_cost += sum(m.c[g]*m.u[g,t]
                         for g in m.G for t in m.T)
    if CONFIG["obj_startup_cost"]:
        total_cost += sum(m.SU_cost[g]*m.y[g,t]
                         for g in m.G for t in m.T)
    if CONFIG["obj_shutdown_cost"]:
        total_cost += sum(m.SD_cost[g]*m.z[g,t]
                         for g in m.G for t in m.T)
    return total_cost
\end{lstlisting}

\subsection{Constraint UC Standar}

\textbf{Keseimbangan daya (Cell 29):}
\begin{equation}
  \sum_{i} P_{i,t} + \sum_{w} P_{\mathrm{WT},w,t}
  + \sum_{b} (P_{d,b,t} - P_{c,b,t}) = D_t \quad \forall\,t
  \label{eq:balance}
\end{equation}

\textbf{Batas kapasitas (Cell 32--33):}
\begin{align}
  u_{i,t}\,P_i^{\min} &\leq P_{i,t} \leq u_{i,t}\,P_i^{\max}
    \quad \forall\,i \in \mathcal{G}_{\mathrm{th}},\,t \label{eq:cap_th}\\
  0 &\leq P_{\mathrm{WT},w,t} \leq P_{\mathrm{WT},w,t}^{\mathrm{avail}}
    \quad \forall\,w,\,t \label{eq:cap_wt}
\end{align}

\textbf{Batas ramp (Cell 38--39):}
\begin{align}
  P_{i,t} - P_{i,t-1} &\leq RU_i\,u_{i,t-1} + P_i^{\min}\,v_{i,t}
    \label{eq:ramp_up}\\
  P_{i,t-1} - P_{i,t} &\leq RD_i\,u_{i,t} + P_i^{\min}\,w_{i,t}
    \label{eq:ramp_down}
\end{align}

\textbf{Waktu nyala/mati minimum (Cell 40--41):}
\begin{align}
  \sum_{\tau=t}^{t+UT_i-1} u_{i,\tau} &\geq UT_i \cdot v_{i,t}
    \label{eq:mut}\\
  \sum_{\tau=t}^{t+DT_i-1} (1-u_{i,\tau}) &\geq DT_i \cdot w_{i,t}
    \label{eq:mdt}
\end{align}
Implementasi menggunakan formulasi \textit{backward-looking}
yang menangani kondisi awal dengan benar.

\textbf{Konsistensi logis (Cell 36--37):}
\begin{align}
  v_{i,t} - w_{i,t} &= u_{i,t} - u_{i,t-1} \label{eq:logic}\\
  v_{i,t} + w_{i,t} &\leq 1 \label{eq:non_sim}
\end{align}

\textbf{N-1 Spinning Reserve (Cell 43):}
\begin{align}
  P_{i,t} + R_{i,t} &\leq P_i^{\max}\,u_{i,t}
    \label{eq:res_headroom}\\
  \sum_{i \in \mathcal{G}_{\mathrm{th}}} R_{i,t}
    &\geq \Ploss(t) \label{eq:n1_res}
\end{align}

\subsection{Constraint BESS}

Selain SoC~\eqref{eq:soc} dan batas daya~\eqref{eq:bess_dis_vi}:
\begin{align}
  \delta_{c,b,t} + \delta_{d,b,t} &\leq 1
    \quad \text{(no simultaneous)} \label{eq:no_sim_cd}\\
  E_b^{\min} \leq E_{b,t} &\leq E_b^{\max}
    \label{eq:soc_bounds}\\
  E_{b,T} &= E_{b,0}
    \quad \text{(energy neutrality)} \label{eq:energy_neutral}
\end{align}

\subsection{Constraint Keamanan Frekuensi}

\subsubsection{Kehilangan Daya Terbesar (Cell 23)}

\begin{equation}
  \Ploss(t) \geq P_{i,t}
    \quad \forall\,i \in \mathcal{G}_{\mathrm{th}},\,t
  \label{eq:largest_loss}
\end{equation}

\begin{lstlisting}[style=pycode,caption={Constraint LargestLoss (Cell 23)}]
def largest_loss_rule(m, g, t):
    if Type[g] in ["PV", "WT"]:
        return Constraint.Skip
    return m.LargestLoss[t] >= m.P[g, t]
model.LargestLossConstraint = Constraint(
    model.G, model.T, rule=largest_loss_rule)
\end{lstlisting}

\subsubsection{Constraint RoCoF (Cell 54 dan 61)}

Dari Persamaan~\eqref{eq:rocof}:
\begin{equation}
  \boxed{
    2\,\RoCoF^{\max}\,\Hsys(t) \geq f_0\,\Ploss(t)
  }
  \quad \forall\,t
  \label{eq:rocof_con}
\end{equation}
Constraint ini \textbf{linier} dalam $\Hsys(t)$ yang
merupakan fungsi linier dari $u_{i,t}$, $P_{\mathrm{VI,BESS},t}$,
dan $P_{\mathrm{VI,WT},t}$.

\begin{lstlisting}[style=pycode,caption={Constraint RoCoF (Cell 54 dan 61)}]
# Cell 54: MinInertiaConstraint
def min_inertia_rule(m, t):
    return (2*RoCoF_limit*m.Hsys[t] >= f0*m.LargestLoss[t])
model.MinInertiaConstraint = Constraint(model.T, rule=min_inertia_rule)

# Cell 61: RoCoFLimitConstraint (ekuivalen matematis)
def rocof_limit_rule(m, t):
    return (m.LargestLoss[t]*f0 <= 2*RoCoF_limit*m.Hsys[t])
model.RoCoFLimitConstraint = Constraint(model.T, rule=rocof_limit_rule)
\end{lstlisting}

\textit{Catatan:} Cell 54 dan 61 mengimplementasikan
constraint yang matematically identik dari perspektif
yang berbeda (inersia minimum vs. batas RoCoF).
Hanya salah satu yang diaktifkan melalui \texttt{CONFIG}.

\subsubsection{Constraint QSS (Cell 62)}

Dari Persamaan~\eqref{eq:qss}:
\begin{equation}
  \Ploss(t) \leq
    (f_0 - f_{\mathrm{QSS}}^{\min})
    \cdot \Big(D\,S_{\mathrm{sys}} + R_{\mathrm{sys},t}\Big)
  \label{eq:qss_con}
\end{equation}

\begin{lstlisting}[style=pycode,caption={Constraint QSS (Cell 62)}]
def qss_limit_rule(m, t):
    damping_term = D_frac * Demand[t]
    return (
        m.LargestLoss[t]
        <= (f0 - f_min) * (damping_term + m.TotalPFR[t])
    )
model.QSSLimitConstraint = Constraint(model.T, rule=qss_limit_rule)
\end{lstlisting}

\subsubsection{Constraint Nadir Frekuensi}

Menggunakan linearisasi dari~\cite{Shi2025}:
\begin{equation}
  \alpha_0 + \alpha_1 \Hsys(t) + \alpha_2 R_{\mathrm{sys},t}
  + \alpha_3 \Ploss(t) \leq \Delta f^{\max}
  \label{eq:nadir_con}
\end{equation}
Dalam kode, \texttt{con\_freq\_nadir} tersedia dalam CONFIG
dan memerlukan prakomputasi koefisien regresi terpisah.

%-----------------------------------------------------------
\section{Konfigurasi Skenario}
\label{sec:config}
%-----------------------------------------------------------

\subsection{Parameter Frekuensi Global (Cell 6)}

\begin{lstlisting}[style=pycode,caption={Parameter frekuensi global}]
f0 = 50.0         # Frekuensi nominal [Hz]
f_min = 49.0      # Batas nadir minimum [Hz]
D_frac = 0.01     # Load damping [%MW/Hz]
RoCoF_limit = 0.55   # Batas RoCoF [Hz/s]
f_qss_min = 49.5     # Batas QSS [Hz]
\end{lstlisting}

\subsection{Flag Konfigurasi Modular}

\begin{table}[!t]\renewcommand{\arraystretch}{1.2}
\caption{Flag Konfigurasi dan Fungsinya}
\label{tab:config}
\centering\footnotesize
\begin{tabular}{lp{5.5cm}}
\toprule
Flag CONFIG & Fungsi \\ \midrule
\texttt{con\_battery}        & Aktifkan model BESS \\
\texttt{con\_battery\_vi}    & Aktifkan VI dari BESS \\
\texttt{con\_wt}             & Aktifkan model turbin angin \\
\texttt{con\_wt\_vi}         & Aktifkan VI dari turbin angin \\
\texttt{con\_min\_inertia}   & Constraint inersia minimum \\
\texttt{con\_rocof\_limit}   & Constraint batas RoCoF \\
\texttt{con\_qss\_limit}     & Constraint batas QSS \\
\texttt{con\_freq\_nadir}    & Constraint nadir (linearisasi) \\
\texttt{con\_n1\_spinning\_reserve} & N-1 reserve konvensional \\
\bottomrule
\end{tabular}
\end{table}

\subsection{Matriks Skenario}

\begin{table}[!t]\renewcommand{\arraystretch}{1.2}
\caption{Matriks Skenario Simulasi}
\label{tab:scenarios}
\centering\footnotesize
\begin{tabular}{lccccc}
\toprule
Skenario & IBR & BESS VI & WT VI & Freq. Con. & Deskripsi \\
\midrule
S0 & 0\%  & -- & -- & -- & UC Klasik (baseline) \\
S1 & 0\%  & -- & -- & Ya & FCUC, inersia SG saja \\
S2 & 40\% & Ya & -- & Ya & FCUC + BESS VI \\
S3 & 40\% & Ya & Ya & Ya & FCUC + BESS + WT VI \\
S4 & 40\% & Ya & Ya & Ya & \textbf{Proposed} (multi-VI) \\
S5 & 60\% & Ya & Ya & Ya & Proposed, 60\% IBR \\
S6 & 80\% & Ya & Ya & Ya & Proposed, 80\% IBR \\
\bottomrule
\end{tabular}
\end{table}

%-----------------------------------------------------------
\section{Parameter Sistem Uji}
\label{sec:params}
%-----------------------------------------------------------

Sistem uji dimodifikasi dari benchmark IEEE 10-unit
Kazarlis et al.~\cite{Kazarlis1996}:

\begin{table}[!t]\renewcommand{\arraystretch}{1.2}
\caption{Parameter Generator Thermal}
\label{tab:gen}
\centering\footnotesize
\begin{tabular}{cccccccc}
\toprule
Unit & $P^{\min}$ & $P^{\max}$ & $b$ & $c$ &
  $H_{\mathrm{SG}}$ & $UT$ & $DT$ \\
 & [MW] & [MW] & [\$/MWh] & [\$/h] & [MWs/MVA] & [h] & [h] \\
\midrule
G1  & 150 & 455 & 16.19 & 1000 & 5.0 & 8 & 8 \\
G2  & 150 & 455 & 17.26 &  970 & 5.0 & 8 & 8 \\
G3  & 20  & 130 & 16.60 &  700 & 4.0 & 5 & 5 \\
G4  & 20  & 130 & 16.50 &  680 & 4.0 & 5 & 5 \\
G5  & 25  & 162 & 19.70 &  450 & 3.5 & 6 & 6 \\
G6  & 20  &  80 & 22.26 &  370 & 3.5 & 3 & 3 \\
G7  & 25  &  85 & 27.74 &  480 & 3.0 & 3 & 3 \\
G8  & 10  &  55 & 25.92 &  660 & 2.5 & 1 & 1 \\
G9  & 10  &  55 & 27.27 &  665 & 2.5 & 1 & 1 \\
G10 & 10  &  55 & 27.79 &  670 & 2.5 & 1 & 1 \\
\bottomrule
\end{tabular}
\end{table}

\begin{table}[!t]\renewcommand{\arraystretch}{1.2}
\caption{Parameter BESS dan IBR}
\label{tab:ibr}
\centering\footnotesize
\begin{tabular}{lrc}
\toprule
Parameter & Nilai & Satuan \\ \midrule
Kapasitas daya BESS & 200 & MW \\
Kapasitas energi BESS & 400 & MWh \\
Batas SoC & 10\%--90\% & -- \\
Efisiensi $\eta_c$, $\eta_d$ & 0.95, 0.95 & -- \\
Biaya degradasi $c_{\mathrm{deg}}$ & 5.0 & USD/MWh \\
Gain VI BESS ($K_{\mathrm{VI}}^{\mathrm{BESS}}$) & 10 & MWs/MW \\
Kapasitas turbin angin & 300 & MW \\
Gain VI WT ($K_{\mathrm{VI}}^{\mathrm{WT}}$) & 3.5 & MWs/MW \\
Kapasitas solar PV & 200 & MW \\
Gain VI PV ($K_{\mathrm{VI}}^{\mathrm{PV}}$) & 2.5 & MWs/MW \\
$S_{\mathrm{sys}}$ & 3000 & MVA \\
$D$ (\textit{load damping}) & 1.0 & \%MW/Hz \\
\bottomrule
\end{tabular}
\end{table}

\begin{table}[!t]\renewcommand{\arraystretch}{1.2}
\caption{Batas Keamanan Frekuensi (sesuai kode)}
\label{tab:freq}
\centering\footnotesize
\begin{tabular}{lrc}
\toprule
Parameter & Nilai & Satuan \\ \midrule
$f_0$ (nominal) & 50.0 & Hz \\
$\RoCoF^{\max}$ & 0.55 & Hz/s \\
$f_{\mathrm{nadir}}^{\min}$ & 49.0 & Hz \\
$f_{\mathrm{QSS}}^{\min}$ & 49.5 & Hz \\
$\Delta f^{\max}$ & 1.0 & Hz \\
\bottomrule
\end{tabular}
\end{table}

%-----------------------------------------------------------
\section{Peta Cell ke Komponen Matematis}
\label{sec:flow}
%-----------------------------------------------------------

\begin{table}[!t]\renewcommand{\arraystretch}{1.2}
\caption{Korespondensi Cell Notebook dengan Persamaan}
\label{tab:cell_map}
\centering\footnotesize
\begin{tabular}{clp{4.8cm}}
\toprule
Cell & Komponen & Persamaan / Deskripsi \\
\midrule
2  & Import       & Library: Pyomo, Pandas, NumPy \\
6  & CONFIG       & Skenario dan parameter frekuensi \\
8  & Data         & Baca Excel: Gen, Demand, Battery, WT \\
12 & Dictionary   & Konversi DataFrame ke dict \\
14 & Derived      & RenewPotential, FreqBias \\
16 & Sets         & model.G, model.T, model.BATT, model.WT \\
18 & Params       & model.Pmin, Pmax, H, Kb\_VI, Kw\_VI \\
20 & Variables    & P, u, y, z, R, P\_VI\_batt, P\_VI\_wt \\
22 & Hsys         & Pers.~\eqref{eq:hsys} \\
23 & LargestLoss  & Pers.~\eqref{eq:largest_loss} \\
24 & TotalPFR     & $R_{\mathrm{sys},t}$ \\
25 & RoCoF\_est   & Pers.~\eqref{eq:rocof} \\
26 & f\_qss       & Pers.~\eqref{eq:qss} \\
28 & Objective    & Pers.~\eqref{eq:obj} \\
29 & Balance      & Pers.~\eqref{eq:balance} \\
32 & MaxCap       & Pers.~\eqref{eq:cap_th}--\eqref{eq:cap_wt} \\
33 & MinCap       & Batas minimum kapasitas \\
36--37 & Transition & Pers.~\eqref{eq:logic}--\eqref{eq:non_sim} \\
38--39 & Ramp      & Pers.~\eqref{eq:ramp_up}--\eqref{eq:ramp_down} \\
40--41 & MUT/MDT   & Pers.~\eqref{eq:mut}--\eqref{eq:mdt} \\
43 & N-1 Reserve  & Pers.~\eqref{eq:res_headroom}--\eqref{eq:n1_res} \\
52 & Battery      & Pers.~\eqref{eq:soc}--\eqref{eq:energy_neutral} \\
54 & MinInertia   & Pers.~\eqref{eq:rocof_con} \\
55 & WT-VI-CL     & Pers.~\eqref{eq:wt_vi_curtailment} \\
61 & RoCoFLimit   & Pers.~\eqref{eq:rocof_con} (ekuivalen) \\
62 & QSSLimit     & Pers.~\eqref{eq:qss_con} \\
64 & Solve        & CPLEX Direct \\
65--83 & Post-proc  & Tabel dan ringkasan \\
87--99 & Freq Sim  & Validasi time-domain \\
101--104 & Plot    & Visualisasi dispatch dan frekuensi \\
105 & Export       & Output ke Excel \\
\bottomrule
\end{tabular}
\end{table}

%-----------------------------------------------------------
\section{Diskusi: Peran Virtual Inertia dalam FCUC}
\label{sec:discussion}
%-----------------------------------------------------------

\subsection{Trade-off Biaya--Keamanan}

Tanpa VI dari IBR, satu-satunya cara memenuhi constraint
RoCoF~\eqref{eq:rocof_con} adalah mengoperasikan cukup
banyak unit thermal. Pada penetrasi IBR tinggi ($>60\%$),
hal ini memaksa unit mahal tetap beroperasi
(\textit{must-run}).

Dengan VI dari IBR, optimizer dapat mengurangi komitmen
unit thermal dan menggantinya dengan alokasi VI:
\begin{equation}
  \Hsys(t) =
    \underbrace{\sum_i H_{\mathrm{SG},i} u_{i,t}}_{\text{dapat dikurangi}}
    + \underbrace{K_{\mathrm{VI}}^{\mathrm{BESS}} P_{\mathrm{VI,BESS}}
      + K_{\mathrm{VI}}^{\mathrm{WT}} P_{\mathrm{VI,WT}}}_{\text{kompensasi VI}}
  \geq H_{\mathrm{req}}(t)
  \label{eq:tradeoff}
\end{equation}

\subsection{Coupling Variabel dalam MILP}

Constraint VI menciptakan \textit{coupling} antara:
\begin{itemize}
  \item Variabel biner $u_{i,t}$ (status unit thermal)
  \item Variabel kontinyu $P_{\mathrm{VI,BESS},t}$,
    $P_{\mathrm{VI,WT},t}$ (alokasi VI)
  \item Variabel kontinyu $P_{d,b,t}$ (dispatch BESS)
  \item Variabel kontinyu $P_{\mathrm{Curt,WT},w,t}$
    (curtailment WT untuk cadangan VI)
\end{itemize}

CPLEX Direct secara simultan mengoptimalkan semua variabel
ini, menemukan kompromi biaya terbaik antara: (a) menambah
unit thermal, (b) mengalokasikan VI dari BESS, dan
(c) mengkurtail WT untuk cadangan VI.

\subsection{Pemetaan Referensi}

\begin{table}[!t]\renewcommand{\arraystretch}{1.2}
\caption{Pemetaan Konsep ke Referensi}
\label{tab:refs}
\centering\footnotesize
\begin{tabular}{lll}
\toprule
Konsep & Referensi Utama & File Lokal \\
\midrule
Inersia \& swing eq. & Kundur~\cite{Kundur1994} & -- \\
Model SFR & Anderson~\cite{Anderson1990} & -- \\
FCUC + BESS VI & Wen et al.~\cite{Wen2016} & \texttt{[01-005]} \\
VI dari turbin angin & Chu et al.~\cite{Chu2020} & \texttt{[01-007]} \\
Linearisasi nadir & Shi et al.~\cite{Shi2025} & \texttt{[01-012]} \\
Review IBR+FCUC & Asiedu et al.~\cite{Asiedu2025} & \texttt{[B09]} \\
Sulawesi UC & \cite{Sulawesi2024} & \texttt{[E02]} \\
JAMALI FFR UC & \cite{JAMALI2025} & \texttt{[E03]} \\
Benchmark 10-unit & Kazarlis et al.~\cite{Kazarlis1996} & -- \\
RUPTL 2025--2034 & PLN~\cite{PLN2025} & \texttt{[RUPTL]} \\
Grid code Indonesia & ESDM~\cite{ESDM2007} & -- \\
\bottomrule
\end{tabular}
\end{table}

%-----------------------------------------------------------
\section{Kesimpulan}
\label{sec:conclusion}
%-----------------------------------------------------------

Dokumen ini telah menyajikan:
\begin{enumerate}
  \item \textbf{Landasan fisik} virtual inertia dari
    perspektif swing equation dan model SFR.
  \item \textbf{Formulasi matematis lengkap} FCUC dengan
    multi-sumber VI (BESS, turbin angin, solar PV).
  \item \textbf{Korespondensi satu-ke-satu} antara
    persamaan matematis dengan implementasi Python/Pyomo
    dalam notebook \texttt{Coba\_16}.
  \item \textbf{Justifikasi fisik} mengapa VI dari IBR
    dapat menggantikan inersia rotasional SG.
  \item \textbf{Pemetaan referensi} menghubungkan setiap
    konsep dengan literatur yang relevan.
\end{enumerate}

Pemodelan ini relevan untuk kondisi sistem tenaga Indonesia,
di mana RUPTL 2025--2034 menargetkan 61\% kapasitas
terbarukan. Pendekatan multi-sumber VI memungkinkan
peningkatan penetrasi IBR sekaligus menjaga keamanan
frekuensi dengan biaya premium minimal.

\bibliographystyle{IEEEtran}
\begin{thebibliography}{20}

\bibitem{PLN2025}
PLN, ``RUPTL PLN 2025--2034,'' Jakarta, 2025.

\bibitem{Tielens2016}
P.~Tielens and D.~Van~Hertem,
``The relevance of inertia in power systems,''
\textit{Renew. Sustain. Energy Rev.}, vol.~55,
pp.~999--1009, 2016.

\bibitem{ESDM2007}
Peraturan Menteri ESDM No.~3 Tahun 2007:
Aturan Jaringan Sistem Tenaga Listrik Jawa-Madura-Bali.

\bibitem{Padhy2004}
N.~Padhy, ``Unit commitment --- a bibliographical survey,''
\textit{IEEE Trans. Power Syst.}, vol.~19, no.~2,
pp.~1196--1205, May 2004.

\bibitem{Wen2016}
Y.~Wen \textit{et al.},
``Frequency dynamics constrained unit commitment with battery
energy storage,''
\textit{IEEE Trans. Power Syst.}, vol.~31, no.~6,
pp.~5115--5125, Nov. 2016. \textit{[Ref 01-005]}

\bibitem{Chu2020}
Z.~Chu \textit{et al.},
``Towards optimal system scheduling with synthetic inertia
provision from wind turbines,''
\textit{IEEE Trans. Power Syst.}, vol.~35, no.~5,
pp.~4056--4066, Sep. 2020. \textit{[Ref 01-007]}

\bibitem{Shi2025}
Q.~Shi \textit{et al.},
``Two-stage linearization of frequency nadir constraint
for unit commitment,''
\textit{IEEE Trans. Power Syst.}, vol.~40, no.~4,
pp.~3584--3587, Jul. 2025. \textit{[Ref 01-012]}

\bibitem{Sulawesi2024}
[Penulis], ``Unit Commitment with Primary Frequency
Regulation in Southern Sulawesi,''
\textit{Proc. ICT-PEP}, 2024. \textit{[Ref E02]}

\bibitem{JAMALI2025}
[Penulis], ``Unit Commitment with 2-Minute FFR Constraint
in JAMALI,''
\textit{Proc. ICT-PEP}, 2025. \textit{[Ref E03]}

\bibitem{Kazarlis1996}
S.~A.~Kazarlis \textit{et al.},
``A genetic algorithm solution to the unit commitment problem,''
\textit{IEEE Trans. Power Syst.}, vol.~11, no.~1,
pp.~83--92, Feb. 1996.

\bibitem{Kundur1994}
P.~Kundur, \textit{Power System Stability and Control}.
New York: McGraw-Hill, 1994.

\bibitem{Anderson1990}
P.~M.~Anderson and M.~Mirheydar,
``A low-order system frequency response model,''
\textit{IEEE Trans. Power Syst.}, vol.~5, no.~3,
pp.~720--729, Aug. 1990.

\bibitem{Asiedu2025}
S.~T.~Asiedu \textit{et al.},
``A review of dynamic constraint integration into
steady-state power system optimization models for
IBR-dominated systems,''
\textit{IEEE Access}, vol.~13, 2025. \textit{[Ref B09]}

\end{thebibliography}

\end{document}